\documentclass[11pt,a4paper]{article}
\usepackage[T1]{fontenc}
\usepackage{lmodern}
\usepackage[margin=25mm]{geometry}
\usepackage{amsmath,amssymb,amsthm,booktabs,microtype}
\usepackage[colorlinks=true,linkcolor=blue,citecolor=blue,urlcolor=blue]{hyperref}
\hypersetup{pdftitle={Four minimal relations in a nearly Gorenstein numerical semigroup},pdfauthor={GPT-6 Astra}}
\newtheorem{theorem}{Theorem}
\newtheorem{lemma}[theorem]{Lemma}
\newcommand{\NN}{\mathbb N}
\newcommand{\ZZ}{\mathbb Z}
\newcommand{\QQ}{\mathbb Q}
\newcommand{\PF}{\operatorname{PF}}
\newcommand{\frakm}{\mathfrak m}
\title{Four minimal relations in a nearly Gorenstein\newline numerical semigroup}
\author{GPT-6 Astra}
\date{November 1, 2026}
\begin{document}
\maketitle
\begin{abstract}
The numerical semigroup $S=\langle6,7,8,9\rangle$ is nearly Gorenstein,
is not almost symmetric, and has exactly four minimal defining relations.
It therefore gives a negative answer to the assertion $\rho(S)\in\{5,6\}$
in Problem 23 of version 3 of the survey by Moscariello and Sammartano.
The example and its four-relation presentation are already in the literature.
We give a self-contained verification of every hypothesis, an elementary
proof of generation and minimality of the defining ideal, and an independent
exact certificate using factorization graphs. We also distinguish the survey's
formulation from the original question of Moscariello and Strazzanti, which
expressly includes four relations in type two.
\end{abstract}

\section{The statement and its sources}
Throughout, $\NN=\{0,1,2,\ldots\}$. A numerical semigroup is an additive
submonoid $S\subseteq\NN$ with finite complement. Its embedding dimension
is the number of minimal positive generators. Write
\[
 F=\max(\ZZ\setminus S),\qquad
 \PF(S)=\{f\in\ZZ\setminus S:f+s\in S\text{ for every }s\in S\setminus\{0\}\}.
\]
The type is $t(S)=|\PF(S)|$. For minimal generators $g_1,\ldots,g_e$, the
semigroup is \emph{nearly Gorenstein} if there are $p_i\in\PF(S)$ such that
\begin{equation}\label{eq:ng}
 g_i+p_i-q\in S\qquad(1\leq i\leq e,\ q\in\PF(S)).
\end{equation}
It is \emph{almost symmetric} if every $a\in\ZZ\setminus S$ satisfies either
$F-a\in S$ or $a,F-a\in\PF(S)$.

Let $I_S$ be the kernel of the homomorphism
$\QQ[x_1,\ldots,x_e]\longrightarrow\QQ[t]$ sending $x_i$ to $t^{g_i}$,
and let $\rho(S)$ be the minimum number of binomials generating $I_S$.
Problem 23 in \cite[version 3]{survey} asks whether
\begin{equation}\label{eq:target}
 e=4,\quad S\text{ nearly Gorenstein and not almost symmetric}
 \quad\Longrightarrow\quad \rho(S)\in\{5,6\}.
\end{equation}

There is a material difference between this formulation and its cited source.
Moscariello and Strazzanti \cite[Question 4.3]{ms} ask whether type two gives
four or five relations, and type three gives six. Their Proposition 3.4 also
places semigroups generated by an arithmetic sequence in the nearly
Gorenstein class. Gomes, O'Neill, and Torres Davila \cite[Examples 8 and 18]{gomes}
give the Ap\'ery set and a four-relation minimal presentation of
$\langle6,7,8,9\rangle$ explicitly. Thus the example below is a known case,
and its relation count belongs to the type-two branch of the original question.
The purpose of this note is to make its consequence for \eqref{eq:target}
explicit and fully verifiable.

\begin{theorem}\label{thm:main}
The semigroup $S=\langle6,7,8,9\rangle$ has embedding dimension four,
$F=11$, and $\PF(S)=\{10,11\}$. It is nearly Gorenstein and is not almost
symmetric. Under $(x,y,z,w)\mapsto(t^6,t^7,t^8,t^9)$, its defining ideal is
\begin{equation}\label{eq:ideal}
 I_S=(y^2-xz,\ yz-xw,\ z^2-yw,\ w^2-x^3),
\end{equation}
and every ideal-generating set has at least four elements. In particular,
$\rho(S)=4$, so \eqref{eq:target} is false.
\end{theorem}

\section{The semigroup hypotheses}
The numbers $6,7,8,9$ are all minimal generators: a sum of two positive
elements is at least $12$. The six consecutive integers from $12$ through
$17$ belong to $S$, since
\[
12=6+6,\quad13=6+7,\quad14=6+8,\quad15=6+9,\quad16=8+8,\quad17=8+9.
\]
Adding $6$ repeatedly proves that every integer at least $12$ belongs to $S$.
Consequently its nonnegative gaps are exactly
\begin{equation}\label{eq:gaps}
 \NN\setminus S=\{1,2,3,4,5,10,11\},\qquad F=11.
\end{equation}

For a gap $f$, testing $f+g_i\in S$ at the minimal generators suffices for
the pseudo-Frobenius condition: every positive element of $S$ is a generator
plus an element of $S$. Both $10$ and $11$ pass this test, since their sums
with generators are at least $16$. Each smaller nonnegative gap fails:
\[
1+9=10,\quad2+8=10,\quad3+7=10,\quad4+6=10,\quad5+6=11.
\]
Negative integers contribute no pseudo-Frobenius numbers. If $f<-6$, then
$f+6<0$. If $-6<f<0$, then $f+6\in\{1,\ldots,5\}$. Finally, $-6+7=1$.
These are all gaps. This proves the asserted full set $\PF(S)=\{10,11\}$.

In the generator order $(6,7,8,9)$, take the NG-vector
$(p_1,p_2,p_3,p_4)=(11,10,10,10)$. All tests in \eqref{eq:ng} are displayed
in Table~\ref{tab:ng}; every entry is a generator of $S$.

\begin{table}[ht]
\centering
\begin{tabular}{rrrr}
\toprule
$g_i$ & $p_i$ & $g_i+p_i-10$ & $g_i+p_i-11$\\
\midrule
6&11&7&6\\
7&10&7&6\\
8&10&8&7\\
9&10&9&8\\
\bottomrule
\end{tabular}
\caption{The complete nearly Gorenstein test.}\label{tab:ng}
\end{table}

Almost symmetry fails at $a=1$: both $1$ and $F-1=10$ are gaps, and
$1\notin\PF(S)$.

\section{Generation and minimality of the defining ideal}
Put $R=\QQ[x,y,z,w]$, and write $\varphi:R\to\QQ[t]$ for the map in
Theorem~\ref{thm:main}. Let $J$ be the ideal on the right of \eqref{eq:ideal}.
Each displayed binomial maps to zero, so $J\subseteq I_S$.

\subsection{A terminating reduction proves generation}
Use the four replacements
\begin{equation}\label{eq:rules}
 y^2\longrightarrow xz,\qquad yz\longrightarrow xw,\qquad
 z^2\longrightarrow yw,\qquad w^2\longrightarrow x^3.
\end{equation}
Give $(x,y,z,w)$ the auxiliary weights $(1,4,6,7)$. These are separate from
their semigroup degrees $(6,7,8,9)$. The four replacements strictly decrease
auxiliary weight, respectively by $1,2,1,11$. Each sequence of replacements
therefore terminates.

A terminal monomial is divisible by none of $y^2,yz,z^2,w^2$. It has the form
\begin{equation}\label{eq:basis}
 x^a b,\qquad a\geq0,\qquad b\in\{1,y,z,w,yw,zw\}.
\end{equation}
The semigroup degrees of the six choices for $b$ are
\[
 0,7,8,9,16,17,
\]
whose residues modulo $6$ are all distinct. If two terminal monomials have
the same image under $\varphi$, reduction of their degrees modulo $6$ first
identifies their choices of $b$ and then identifies their exponents $a$.
Thus all terminal monomials have distinct images, which are linearly
independent powers of $t$.

For any $P\in I_S$, reduce its monomials using \eqref{eq:rules}, collecting
terms at the end, to obtain a linear combination $N$ of terminal monomials.
Every replacement changes $P$ by an element of $J$, so $P-N\in J$ and
$\varphi(N)=0$. The independence just proved gives $N=0$. Hence $P\in J$,
proving $I_S=J$ and $\rho(S)\leq4$.

\subsection{Four relations are necessary}
Use the positive grading $\deg(x,y,z,w)=(6,7,8,9)$. The complete list of
monomials of degree less than $14$ is
\[
 1,x,y,z,w,x^2,xy,
\]
with distinct degrees $0,6,7,8,9,12,13$. Therefore $I_S$ has no nonzero
homogeneous part below degree $14$. For $\frakm=(x,y,z,w)$, the ideal
$\frakm I_S$ has no nonzero homogeneous part below degree $20=6+14$.

The four binomials in \eqref{eq:ideal} have distinct degrees $14,15,16,18$,
all below $20$. Their classes in the graded $\QQ$-vector space
$I_S/\frakm I_S$ are nonzero and linearly independent. Any ideal-generating
set spans this quotient, since its polynomial coefficients reduce to constants
modulo $\frakm$. Every such set therefore has at least four elements.
Together with the generation proof, this proves $\rho(S)=4$ and completes
Theorem~\ref{thm:main}.

\section{An independent exact check by factorization graphs}
For $n\in S$, let
\[
 Z(n)=\{(a,b,c,d)\in\NN^4:6a+7b+8c+9d=n\}.
\]
Join two distinct factorizations by an edge when they share a positive
coordinate, and denote the number of connected components by $c(n)$.
The number of minimal relations of degree $n$ is $c(n)-1$; this is the
standard presentation criterion recalled in \cite[Theorem 3]{gomes}.
Here is the algebraic reason. The degree-$n$ part of $I_S$ consists of
linear combinations of monomials in $Z(n)$ whose coefficients sum to zero.
The degree-$n$ part of $\frakm I_S$ is spanned by differences of monomials
sharing a variable, hence by differences within graph components.
Its quotient therefore has dimension $c(n)-1$.

\begin{lemma}\label{lem:bound}
For a numerical semigroup with four ordered minimal generators
$g_1<g_2<g_3<g_4$ and Frobenius number $F$, every disconnected
factorization graph has degree at most $F+g_3+g_4$.
\end{lemma}
\begin{proof}
If $n>F+g_3+g_4$, then $n-g_i-g_j>F$ for every distinct pair $i,j$.
Thus there is a factorization of $n$ containing both generators $g_i,g_j$.
Given any two factorizations, choose a generator from each support.
If the chosen generators coincide, the factorizations are adjacent;
otherwise a factorization containing both connects them by a path of
length at most two. The graph is connected.
\end{proof}

For the present example the bound is $11+8+9=28$. Exact enumeration of
all factorization graphs through degree $28$ finds precisely the four
disconnected graphs in Table~\ref{tab:betti}. All four have two components,
so the independent count is $\rho(S)=1+1+1+1=4$.

\begin{table}[ht]
\centering
\begin{tabular}{rllr}
\toprule
Degree & First factorization & Second factorization & Components\\
\midrule
14 & $y^2$ & $xz$ & 2\\
15 & $yz$ & $xw$ & 2\\
16 & $z^2$ & $yw$ & 2\\
18 & $w^2$ & $x^3$ & 2\\
\bottomrule
\end{tabular}
\caption{All disconnected factorization graphs; each listed fiber has exactly
the two shown monomials.}\label{tab:betti}
\end{table}

The Ap\'ery set, consisting of the smallest elements in each residue class
modulo $6$, is $\{0,7,8,9,16,17\}$. It agrees with the degrees of the
six terminal basis monomials in \eqref{eq:basis} and with
\cite[Example 8]{gomes}.

\section{Reproducibility and attribution}
The accompanying package contains two standard-library Python programs.
From the directory containing \texttt{src/} and \texttt{certificates/}, run
\begin{verbatim}
python src/proof_verify.py
python src/search_verify.py
\end{verbatim}
The first checks the finite data of the terminating-reduction and minimality
arguments. The second independently enumerates all factorizations through
the proved bound, checks the semigroup hypotheses, and compares the complete
result with the supplied JSON certificate. Both use exact integer arithmetic.
The preceding proofs supply the universal steps beyond these finite checks.

This note and its programs were prepared by GPT-6 Astra. The result is a
self-contained verification of a known example and its negative answer to
the precise formulation \eqref{eq:target}. The source comparison and the
literature search are documented in the accompanying literature review.

\begin{thebibliography}{9}
\bibitem{survey}
A. Moscariello and A. Sammartano,
\emph{Open problems on relations of numerical semigroups},
arXiv:2406.00790v3, September 28, 2026, Problem 23.
\url{https://arxiv.org/abs/2406.00790v3}.

\bibitem{ms}
A. Moscariello and F. Strazzanti,
\emph{Nearly Gorenstein vs Almost Gorenstein Affine Monomial Curves},
Mediterranean Journal of Mathematics \textbf{18} (2021), article 127,
Proposition 3.4 and Question 4.3.
\href{https://doi.org/10.1007/s00009-021-01761-1}{doi:10.1007/s00009-021-01761-1}.

\bibitem{gomes}
T. Gomes, C. O'Neill, and E. Torres Davila,
\emph{Numerical Semigroups, Polyhedra, and Posets III:
Minimal Presentations and Face Dimension},
Electronic Journal of Combinatorics \textbf{30}(2) (2023), \#P2.57,
Theorem 3 and Examples 8 and 18.
\href{https://doi.org/10.37236/10380}{doi:10.37236/10380}.
\end{thebibliography}
\end{document}
